\chapter{Resultados e discussão}

A \autoref{tab:criterios} relaciona cada critério de avaliação do enunciado com o
resultado obtido e a evidência correspondente.

\begin{table}[htbp]
  \centering
  \caption{Critérios de avaliação e resultados}\label{tab:criterios}
  \begin{tabular}{L{0.3\textwidth}L{0.38\textwidth}L{0.2\textwidth}}
    \hline
    \textbf{Critério} & \textbf{Resultado} & \textbf{Evidência} \\
    \hline
    Cálculo de sub-redes e configuração de IP & 4 sub-redes /26 de 62 \textit{hosts}; VMs na sub-rede A; cliente na sub-rede B isolado do servidor & \autoref{sec:enderecamento} \\
    Efetividade do \textit{firewall} & antes: 22/tcp e 8080/tcp abertas; depois: somente 8080/tcp (65.534 portas filtradas) & \autoref{sec:firewall} \\
    Menor privilégio no IAM & \texttt{lab-user} lê, lista e grava só no \texttt{lab-bucket}; as 6 ações testadas fora do escopo foram negadas & \autoref{sec:iam}, \autoref{sec:verificacao} \\
    Armazenamento privado & acesso anônimo (leitura, listagem e escrita) passou de HTTP 200 para HTTP 403 & \autoref{sec:iam} \\
    Evidência completa da Fase 4 & 18 testes (13 no servidor e 5 no cliente), nenhuma falha & \autoref{sec:verificacao} \\
    \hline
  \end{tabular}
  \legend{Fonte: elaborado pelo autor.}
\end{table}

\section{Defesa em profundidade}

O laboratório mostra dois controles independentes protegendo o mesmo dado. Mesmo com o
\textit{bucket} público, a API S3 do SILO não era alcançável pelo cliente, porque o
\textit{firewall} só libera a porta 8080. Por outro lado, se a regra de \textit{firewall}
fosse relaxada, a política do \textit{bucket} e do usuário continuariam negando o acesso
indevido. Filtragem de rede responde \emph{quem alcança o serviço}; IAM responde
\emph{quem pode fazer o quê em qual recurso}. Um não substitui o outro.

\section{Equivalência com provedores de nuvem}

Os conceitos exercitados têm correspondência direta nos provedores públicos, o que
permite transpor o laboratório para uma conta real (\autoref{tab:nuvem})
\cite{nist800145,aws-s3-bpa}.

\begin{table}[htbp]
  \centering
  \caption{Correspondência entre o laboratório e a AWS}\label{tab:nuvem}
  \begin{tabular}{L{0.42\textwidth}L{0.48\textwidth}}
    \hline
    \textbf{Laboratório} & \textbf{AWS} \\
    \hline
    Sub-redes /26 em 192.168.50.0/24 & VPC com sub-redes \\
    \texttt{ufw} no servidor & \textit{Security Group} / \textit{Network ACL} \\
    Usuário \texttt{lab-user} + política \texttt{lab-rw} & Usuário ou \textit{role} do IAM + política gerenciada \\
    \texttt{mcli anonymous set public/private} & \textit{Bucket policy} com \texttt{Principal "*"} e \textit{S3 Block Public Access} \\
    \hline
  \end{tabular}
  \legend{Fonte: elaborado pelo autor.}
\end{table}

\section{Limitações}

\begin{alineas}
  \item a política de saída do \texttt{ufw} foi mantida como ``permitir'', pois o enunciado
    trata apenas do tráfego de entrada; em produção, a saída também deveria ser restrita;
  \item a exceção de gerência do Vagrant (porta 22 pela NAT, a partir do hospedeiro) existe
    apenas por causa da ferramenta de automação e não é alcançável pela rede do laboratório;
  \item os testes de acesso anônimo ao \textit{bucket} foram feitos a partir do próprio
    servidor, porque a rede bloqueia a porta 9000; isso isola o teste do controle de IAM
    do teste do controle de rede;
  \item o SILO é um \textit{fork} comunitário; em uma conta de nuvem real, o mesmo roteiro
    usaria IAM e S3 nativos do provedor.
\end{alineas}

\section{Observações da execução}

A varredura completa das 65.535 portas levou cerca de 26 segundos antes do
\textit{firewall} e cerca de 86 segundos depois. A diferença se explica pelo
comportamento de cada estado: uma porta \emph{fechada} responde imediatamente com um
pacote de recusa (\texttt{RST}), enquanto uma porta \emph{filtrada} não responde, e o
\texttt{nmap} precisa esperar o tempo limite de cada tentativa. Do ponto de vista de quem
varre, o \textit{firewall} portanto não só esconde os serviços como encarece o
reconhecimento da rede. O mesmo efeito aparece nos testes de aplicação: a conexão à
porta 22 expira por \textit{timeout}, em vez de ser recusada.
